\section{Exact polynomial certificates}
\label{sec:polynomial-certificates}

We now prove the polynomial inequalities in
Proposition~\ref{prop:polynomial-inequalities}. The single-posterior
inequalities follow from small homogeneous coefficient calculations. For the
two-posterior inequalities, we give two sufficient tests on a product of
simplices. Section~\ref{sec:finite-verification} applies those tests to a
finite subdivision and completes the proof.

\subsection{Ordered probabilities and homogeneous coefficients}

Let
\[
 \mathcal S=\{p\in\Deltafour:p_1\ge p_2\ge p_3\ge p_4\},
 \qquad
 S_i=\frac{1}{i+1}\sum_{v=1}^{i+1}e_v
 \quad(0\le i\le3),
\]
where $e_v$ is the $v$th standard basis vector of $\R^4$.
The set $\mathcal S$ is the simplex with vertices $S_0,S_1,S_2,S_3$.
Indeed, on setting $p_5=0$, its barycentric coordinates are
\begin{equation}
 X_i=(i+1)(p_{i+1}-p_{i+2}),\qquad
 p=\sum_{i=0}^3X_iS_i,\qquad \sum_{i=0}^3X_i=1.
 \label{eq:ordered-barycentric}
\end{equation}
All $X_i$ are nonnegative. Conversely, each such convex combination lies
in $\mathcal S$.

A permutation of the four probabilities acts on $x=Lp$ by permuting its
three coordinates and changing an even number of their signs. One way to
check this is to note that, on the hyperplane $\sum_vp_v=1$,
$p=\unif+L^{\mathsf T}x/4$. For a permutation matrix $\Pi$, the induced
matrix is $L\Pi L^{\mathsf T}/4$; it is a signed permutation matrix whose
three nonzero entries have product $1$. Consequently $A,B,C$, and hence
$H,F,G$, are invariant. It is therefore enough to check the
single-posterior inequalities on $\mathcal S$. The permutation-averaged
two-posterior expressions are invariant under separate permutations of
$p$ and $q$: a permutation of $q$ reindexes the average, and a permutation
of $p$ can be moved to $q$ by simultaneous permutation of the output
coordinates. We may thus restrict both inputs to $\mathcal S$.

We use ordinary monomial coefficients throughout. For $r\ge0$, put
\[
 \mathcal I_r=\{I\in\mathbb Z_{\ge0}^4:|I|=r\},\qquad
 X^I=\prod_{i=0}^3X_i^{I_i},\qquad
 W_r(X)=\left(\sum_{i=0}^3X_i\right)^r.
\]
The coefficient of $X^I$ in $W_r$ is the multinomial coefficient
$\binom rI=r!/\prod_iI_i!$. A polynomial of degree at most $r$ can be
homogenized to degree $r$ by multiplying each degree-$s$ monomial by
$W_{r-s}$. This does not change its values when $\sum_iX_i=1$.
If all coefficients of a homogeneous polynomial are nonnegative, then
the polynomial is nonnegative on the simplex. No converse is required.
The coefficients in the simplicial Bernstein basis differ from these
ordinary homogeneous coefficients by positive multinomial factors, so
their signs agree. This is the coefficient-positivity principle behind
simplicial Bernstein bounds and their refinement by subdivision;
see \cite{Leroy2012}.

The calculations below use only coefficient addition and the product rule
\begin{equation}
 [X^I](RS)=\sum_{J+K=I}[X^J]R\,[X^K]S.
 \label{eq:coefficient-convolution}
\end{equation}
For two sets of variables, write
$W_r(X,Y)=W_r(X)W_r(Y)$. A bihomogeneous polynomial of bidegree $(r,r)$,
meaning homogeneous of degree $r$ in each set of variables, is stored as
the matrix of its coefficients $[X^IY^J]$, with $I,J\in\mathcal I_r$.
Multiplication is convolution in both indices; multiplication by
$W_s(X,Y)$ raises both homogeneous degrees by $s$.
A monomial of degree $s$ in $X$ and degree $t$ in $Y$ is homogenized
to bidegree $(r,r)$ by multiplying it by
$W_{r-s}(X)W_{r-t}(Y)$.

\subsection{The single-posterior checks}

The coordinate forms on $\mathcal S$ are particularly small:
\begin{equation}
 L\sum_{i=0}^3X_iS_i
 =\left(X_0+X_1+\frac{X_2}{3},\,
         X_0+\frac{X_2}{3},\,
         X_0-\frac{X_2}{3}\right).
 \label{eq:ordered-coordinate-forms}
\end{equation}
Let $\Phi_0,\ldots,\Phi_8$ be the polynomials
$1,A,B,A^2,C,AB,A^3,AC,B^2$ in the three coordinates, and let
\[
 (e_0,\ldots,e_8)=(0,2,3,4,4,5,6,6,6)
\]
be their degrees. If $R=\sum_jc_j\Phi_j$, its degree-$r$ homogeneous
form on $\mathcal S$ is
\begin{equation}
 \widehat R_r(X)
 =\sum_{j:c_j\ne0}c_j\,
   \Phi_j\left(L\sum_{i=0}^3X_iS_i\right)W_{r-e_j}(X),
 \qquad r\ge\max_{c_j\ne0}e_j.
 \label{eq:scalar-homogenization}
\end{equation}
Together with \eqref{eq:ordered-coordinate-forms} and
\eqref{eq:coefficient-convolution}, this formula specifies every
coefficient in the following check.

\begin{lemma}[Single-posterior certificate]
\label{lem:scalar-checks}
For every $p\in\Deltafour$,
\[
 H(p)\ge-\frac12,\qquad F(p)\ge0,\qquad G(p)\ge0.
\]
Moreover, if $x=Lp$ and $0\le\ell\le1$, then
$H(\ell x)\ge\ell^3H(x)$, where $H$ in this formula is regarded as a
polynomial in the three coordinates.
\end{lemma}

\begin{proof}
Define the Euler operator on those coordinates by
$\mathcal E=\sum_{j=1}^3x_j\partial_{x_j}$. Thus
$\mathcal E\Phi_j=e_j\Phi_j$. The exact coefficient minima in
Table~\ref{tab:scalar-coefficients} are obtained from
\eqref{eq:scalar-homogenization}. The first row adds $1/2$ to the
constant coefficient of $H$; the last row replaces each coefficient
$c_j$ of $H$ by $(3-e_j)c_j$. The full scalar verifier is included in
Appendix~\ref{app:verifiers}.

\begin{table}[tb]
 \centering
 \begin{tabular}{lrrr}
  \hline
  Polynomial $R$ & Degree $r$ & Number of coefficients & Minimum\\
  \hline
  $H+1/2$          & 6  & 84  & $19/50$\\
  $F$              & 5  & 56  & $0$\\
  $G$              & 14 & 680 & $87/1000$\\
  $(3-\mathcal E)H$ & 6  & 84  & $0$\\
  \hline
 \end{tabular}
 \caption{Minimum ordinary monomial coefficients of the homogeneous
 forms \eqref{eq:scalar-homogenization}, computed exactly over $\mathbb Q$.}
 \label{tab:scalar-coefficients}
\end{table}

Coefficient nonnegativity proves the first three assertions and
$(3-\mathcal E)H\ge0$ on $L\mathcal S$. Permutation invariance extends
these statements to $L\Deltafour$. For $x\in L\Deltafour$ and
$0<\ell\le1$, the point $\ell x$ also lies in $L\Deltafour$, since
it corresponds to $\ell p+(1-\ell)\unif$. Therefore
\[
 \frac{d}{d\ell}\bigl(\ell^{-3}H(\ell x)\bigr)
 =\ell^{-4}(\mathcal E-3)H(\ell x)\le0.
\]
Comparing $\ell$ with $1$ gives the asserted scaling inequality.
At $\ell=0$ it follows from $H(0)=0$.
\end{proof}

\subsection{The two rational inequalities}

It remains to prove the two-posterior inequalities. We express each
as the nonnegativity of a rational function on
$\mathcal S\times\mathcal S$. For $\pi\in\Sym$, let
\[
 U^\pi_v=4p_vq_{\pi(v)},\qquad
 z_\pi=\sum_{v=1}^4U^\pi_v,\qquad
 P^\pi=U^\pi/z_\pi\quad(z_\pi>0).
\]
Set
\begin{align}
 K_a&=24\bigl(A(p)+A(q)-H(p)F(q)-F(p)H(q)-D_0(p,q)\bigr),\notag\\
 K_b&=-24\bigl(H(p)G(q)+G(p)H(q)\bigr).
 \label{eq:rational-K}
\end{align}
The remaining data are
\begin{equation}
 \begin{array}{c|c|c|c}
  \text{row} & k & N_\pi & m\\ \hline
  a & 1 & -\displaystyle\sum_{j=1}^3(LU^\pi)_j^2 & 3\\[2mm]
  b & 5 & z_\pi^6H(U^\pi/z_\pi) & 1/2
 \end{array}
 \label{eq:rational-rows}
\end{equation}
In row $b$, the expression for $N_\pi$ is a polynomial: each
degree-$s$ term of $H$ becomes the corresponding term in $LU^\pi$
multiplied by $z_\pi^{6-s}$. For either row define
\begin{equation}
 \mathcal R=K+\sum_{\pi\in\Sym}\frac{N_\pi}{z_\pi^k}.
 \label{eq:rational-target}
\end{equation}
In row $a$, the summand is $-z_\pi A(P^\pi)$; in row $b$, it is
$z_\pi H(P^\pi)$. Thus $\mathcal R\ge0$ is exactly the corresponding
inequality of Proposition~\ref{prop:polynomial-inequalities}, multiplied
by $24$.

These rational functions have continuous extensions to their apparent
singularities. All coordinates of $U^\pi$ are nonnegative, so
$z_\pi=0$ implies $U^\pi=0$. Since $P^\pi\in\Deltafour$ whenever
$z_\pi>0$, both $A(P^\pi)$ and $H(P^\pi)$ are bounded independently
of $p,q$. Hence the corresponding weighted summand tends to $0$ as
$z_\pi\to0$. We assign it that value on the boundary.

Now take any two nondegenerate subsimplices of $\mathcal S$, with
vertex lists $a=(a_0,\ldots,a_3)$ and $b=(b_0,\ldots,b_3)$, and write
\[
 p=\sum_{i=0}^3X_ia_i,\qquad q=\sum_{j=0}^3Y_jb_j,
 \qquad X_i,Y_j\ge0,\quad \sum_iX_i=\sum_jY_j=1.
\]
The ordinary bidegree-$(1,1)$ coefficient of $U^\pi_v$ at $X_iY_j$
is $4a_{iv}b_{j,\pi(v)}$. In particular, the coefficients of $z_\pi$
are nonnegative. Both $p$ and $q$ are strictly positive on the
relative interior of their subsimplices: a full-dimensional simplex
cannot lie in a coordinate face of $\Deltafour$. Thus every $z_\pi$
is strictly positive there. We represent $K$ homogeneously with
bidegree $(6,6)$ and $N_\pi$ with bidegree $(k+1,k+1)$.

The next two tests certify $\mathcal R\ge0$ on such a product of
simplices. They establish lower bounds on the entire region, not just
at the points used to construct the polynomials.

\subsection{A Taylor lower bound: Type T}

For a positive integer $k$, define
\[
 t_k(z)=\sum_{i=0}^5\binom{-k}{i}(z-1)^i.
\]
For every $z>0$, Taylor's theorem gives
\begin{equation}
 z^{-k}-t_k(z)
 =\frac{k(k+1)\cdots(k+5)}{6!}\,
   \xi^{-k-6}(z-1)^6\ge0
 \label{eq:taylor-lower-bound}
\end{equation}
for some $\xi$ between $1$ and $z$; at $z=1$ the identity is
immediate. The bound holds for the full positive half-line, with no
restriction such as $|z-1|<1$.

For row $a$, $A(P)\le3$ on $\Deltafour$, because each coordinate of
$LP$ lies in $[-1,1]$. For row $b$, Lemma~\ref{lem:scalar-checks}
gives $H(P)\ge-1/2$. Consequently, in either row,
\begin{equation}
 N_\pi+mz_\pi^{k+1}\ge0.
 \label{eq:taylor-nonnegative-factor}
\end{equation}

\begin{lemma}[Type T certificate]
\label{lem:type-t}
For either row of \eqref{eq:rational-rows}, form the polynomial
\begin{equation}
 T=K+\sum_{\pi\in\Sym}
       \bigl((N_\pi+mz_\pi^{k+1})t_k(z_\pi)-mz_\pi\bigr).
 \label{eq:type-t-polynomial}
\end{equation}
If its homogeneous coefficient matrix in bidegree $(k+6,k+6)$
has no negative entries, then $\mathcal R\ge0$ on the product of
simplices.
\end{lemma}

\begin{proof}
On the relative interior, multiply \eqref{eq:taylor-lower-bound} by
the nonnegative factor in \eqref{eq:taylor-nonnegative-factor} and
subtract $mz_\pi$. This yields
\[
 \frac{N_\pi}{z_\pi^k}
 \ge (N_\pi+mz_\pi^{k+1})t_k(z_\pi)-mz_\pi.
\]
Summing gives $\mathcal R\ge T$. The coefficient hypothesis makes
$T$ nonnegative on the product simplex. The continuous extension of
$\mathcal R$ then proves the boundary cases as well.
\end{proof}

\subsection{A polynomial lower bound with a certified remainder: Type P}

The second test constructs a lower bound for each rational summand
separately. Start with any bihomogeneous polynomial $J_\pi$ of bidegree $(3,3)$.
The way it is chosen affects the strength of the test, but not its
validity. We seek a constant $c_\pi\ge0$ such that
\[
 \frac{N_\pi}{z_\pi^k}\ge J_\pi-c_\pi
\]
throughout the product simplex. On the relative interior, clearing the positive
denominator reduces this to nonnegativity of
$N_\pi-z_\pi^kJ_\pi+c_\pi z_\pi^k$.
To test that remainder by homogeneous coefficients, form in bidegree
$(k+3,k+3)$ the arrays
\begin{equation}
 O_\pi=W_2N_\pi-z_\pi^kJ_\pi,\qquad
 Q_\pi=z_\pi^kW_3,
 \label{eq:type-p-remainder}
\end{equation}
where the arguments $(X,Y)$ of $W_2,W_3$ are suppressed.
Every coefficient of $Q_\pi$ is nonnegative.

Fix $M=2^{32}$. If an entry $(O_\pi)_{IJ}$ is negative and the
corresponding entry $(Q_\pi)_{IJ}$ is zero, the test is not
applicable. Otherwise define
\begin{equation}
 c_\pi=\frac1M\max\left(
  \{0\}\cup
  \left\{
   \left\lceil\frac{-M(O_\pi)_{IJ}}{(Q_\pi)_{IJ}}\right\rceil:
   (O_\pi)_{IJ}<0
  \right\}\right).
 \label{eq:remainder-correction}
\end{equation}
By construction, $O_\pi+c_\pi Q_\pi$ is coefficientwise
nonnegative.

\begin{lemma}[Type P certificate]
\label{lem:type-p}
Suppose \eqref{eq:remainder-correction} is defined for every
$\pi\in\Sym$. If
\begin{equation}
 K+\sum_{\pi\in\Sym}(J_\pi-c_\pi W_3)
 \label{eq:type-p-polynomial}
\end{equation}
has nonnegative homogeneous coefficients in bidegree $(6,6)$,
then $\mathcal R\ge0$ on the product of simplices.
\end{lemma}

\begin{proof}
On the product simplex, $W_2=W_3=1$. On its relative interior,
\eqref{eq:type-p-remainder} gives the exact identity
\[
 \frac{N_\pi}{z_\pi^k}
 =J_\pi-c_\pi+
   \frac{O_\pi+c_\pi Q_\pi}{z_\pi^k}
 \ge J_\pi-c_\pi.
\]
Summing and using the coefficient hypothesis proves the claim on the
relative interior. Continuity gives the boundary cases.
\end{proof}

For reproducibility, we specify the polynomials $J_\pi$ used in the
finite check. For $I\in\mathcal I_3$, define the cardinal polynomial
\[
 \ell_I(X)=\prod_{i=0}^3\binom{3X_i}{I_i},\qquad
 \binom{t}{j}=\frac{t(t-1)\cdots(t-j+1)}{j!},\qquad
 \binom{t}{0}=1.
\]
At a lattice point $J/3$, with $J\in\mathcal I_3$, its value is
$\prod_i\binom{J_i}{I_i}=\mathbf 1_{\{I=J\}}$. Indeed, if $I\ne J$
then some $J_i<I_i$, since $|I|=|J|=3$. These $20$ cardinal
polynomials therefore give interpolation of degree at most three
on the three-dimensional simplex. In coefficient arrays we use
their homogeneous degree-three forms, obtained by replacing each
factor $3X_i-h$ by $3X_i-h\sum_jX_j$.

Let $\operatorname{trunc}_M(t)$ denote truncation toward zero to a
multiple of $1/M$. Evaluate the continuous rational summand
$N_\pi/z_\pi^k$ at the $20^2$ node pairs $(I/3,J/3)$, assigning zero
when $z_\pi=0$, and put
\begin{equation}
 J_\pi(X,Y)=
 \sum_{I,J\in\mathcal I_3}
 \operatorname{trunc}_M\!\left(
   \frac{N_\pi(I/3,J/3)}{z_\pi(I/3,J/3)^k}\right)
 \ell_I(X)\ell_J(Y).
 \label{eq:interpolation-polynomial}
\end{equation}
The convention at vanishing denominators is understood in this
formula. For the rational vertex lists used in
Section~\ref{sec:finite-verification}, all these node values are rational.
Their evaluation, truncation, and interpolation can therefore be performed
in exact rational arithmetic. Lemma~\ref{lem:type-p} does not estimate
interpolation error or truncation error: it certifies the remainder
of the actual rational polynomial $J_\pi$ chosen by
\eqref{eq:interpolation-polynomial}. Thus the verification remains
exact despite the finite interpolation grid and truncation.

The single-posterior inequalities are now proved. For the two
remaining inequalities, the proof has been reduced to finding a
finite cover of $\mathcal S\times\mathcal S$ by products of
simplices on each of which either Lemma~\ref{lem:type-t} or
Lemma~\ref{lem:type-p} applies. We carry out that finite verification
next.
